\documentclass[11pt,letterpaper]{article}

\usepackage[utf8]{inputenc}
\usepackage[T1]{fontenc}
\usepackage[left=0.7in,right=0.7in,top=0.6in,bottom=0.6in]{geometry}
\usepackage{enumitem}
\usepackage{titlesec}
\usepackage{hyperref}
\usepackage{xcolor}
\usepackage{graphicx}

\usepackage{microtype}
\usepackage[none]{hyphenat}
\sloppy
\frenchspacing

% Match font and spacing style from CV (no indent, standard Computer Modern, fixed gap)
\setlength{\parindent}{0pt}
\setlength{\parskip}{1ex}

\hypersetup{
    colorlinks=true,
    linkcolor=black,
    filecolor=black,      
    urlcolor=blue,
    citecolor=black,
}

\pagestyle{empty}

\begin{document}

\noindent
% \parbox[b]{0.5\textwidth}{}%
% \parbox[b]{0.5\textwidth}{\raggedright Hanoi, July 20, 2026}

\vspace{3ex}

\begin{center}
    \Large \textbf{STATEMENT OF PURPOSE}
\end{center}

\vspace{2ex}

\noindent \textbf{Applicant:} Xuan Thanh Trinh\\
\noindent \textbf{Program:} Ph.D. in Electrical, Electronics and Computer Engineering\\
\noindent \textbf{Proposed Advisor:} Assistant Professor Van-Hai Bui\\
\noindent \textbf{Institution:} University of Michigan-Dearborn

\vspace{3ex}

\noindent \textbf{1. Introduction}

As a final-year Electrical Engineering student at Hanoi University of Science and Technology (HUST), my research at the Power Grid and Renewable Energy (PGRE) Laboratory focuses on developing distributed mathematical optimization models for decentralized energy management. While this foundation drove my graduation thesis on microgrid resilience, I recognized that classical optimization alone cannot scale to the real-time computational demands of modern power systems. To overcome these computational bottlenecks, I require advanced training in scalable, cyber-physical control systems. The Ph.D. program at UM-Dearborn perfectly aligns with this need, providing the rigorous framework I need to become a highly specialized research engineer capable of designing decentralized architectures for high renewable penetrations.

\vspace{1ex}
\noindent \textbf{2. Academic Background}

Ranking top 3 in HUST's Advanced Program in Electrical Engineering and Renewable Energy instilled in me a systems-level perspective. Rather than isolating individual grid components, I trained to deconstruct complex power networks into dynamic, tractable optimization problems. This mathematical foundation drives my independent research trajectory, allowing me to architect solutions that bridge abstract algorithms with real world physical constraints.

This trajectory reflects a systematic transition from treating power grids as aggregate generation-load balances to modeling them as complex computational graphs. In research recently submitted to \textit{Sustainable Energy, Grids and Networks} (SEGAN), I established and validated a peer-to-peer trading framework that subdivides large-scale distribution networks into autonomous microgrids, integrating rigorous grid physics (AC-OPF) with distributed optimization (ADMM) across both radial and ring-connected topologies. By accounting for marginal transmission losses, this architecture reduced operational costs while strictly bounding voltage levels within the $\pm 5\%$ threshold, preventing local trading from destabilizing the wider public grid. However, this static framework lacked the temporal foresight to handle sudden local failures.

True grid resilience demands dynamic fault tolerance. To address the static model's vulnerability, I dedicated my graduation thesis to layering Model Predictive Control (MPC) over the trading baseline to coordinate emergency power exchanges and backup diesel generation. During simulated localized outages, this temporal coordination empowered neighboring networks to assist a failing grid, reducing forced load shedding from 15 MWh down to just 3 MWh while protecting 100\% of critical loads.

Ultimately, this intensive research trajectory has equipped me with a robust technical toolkit (Python, Pyomo) and a collaborative, resilient mindset. I am ready to bring this research foundation to your university, eager to fuse my expertise with your computational ecosystem to drive scalable, real-time grid operations.

\vspace{1ex}
\noindent \textbf{3. Research Interests}

To resolve the trade-off between computational tractability and physical grid security, my Ph.D. research will focus on architecting Physics-Informed AI frameworks for real-time power system dispatch. While pure data-driven models accelerate computational speeds, their deployment frequently fails because unconstrained black-box policies inherently violate non-linear AC power flow equations and thermal limits. I intend to synthesize these paradigms by deploying my deterministic optimization models in a dual capacity: as expert benchmarks to generate mathematically rigorous optimal dispatch datasets, and as differentiable safety filters—leveraging convex relaxations like SOCP—that project AI-generated setpoints into a verifiable feasible region. By incorporating structural grid constraints into both the offline training and online execution agents, I aim to pioneer certifiably safe and physically resilient dispatch architectures.

\vspace{1ex}
\noindent \textbf{4. Academic and Research Alignment}

Realizing that architecting Physics-Informed AI frameworks requires an ecosystem that seamlessly bridges advanced machine learning with rigorous mathematical optimization, I found the Ph.D. program at UM-Dearborn to be the perfect catalyst for my research. To construct mathematically guaranteed safety filters for these AI agents, I require a stringent theoretical foundation. The interdisciplinary scope of the EECE curriculum will systematically bridge my theoretical gaps, providing the comprehensive mathematical and computational toolkit necessary for this architecture through advanced courses such as IMSE 505 (Optimization) and ECE 5001 (Analytic and Computational Math). Furthermore, the robust computational facilities at UM-Dearborn, paired with its strategic ties to regional energy sectors, provide the exact environment I need to rigorously test and validate these theoretical architectures against real-world grid scenarios.

Within this dynamic ecosystem, my research trajectory strongly aligns with the ongoing work of Prof. Van-Hai Bui in autonomous energy markets. His recent publications on FairMarket-RL and XRL-LLM investigate the integration of Large Language Models into multi-agent energy trading, an area highly relevant to my interest in learning-based policies. By joining his laboratory, my foundation in mathematical optimization will allow me to seamlessly support his ongoing research while advancing my own theoretical frameworks. Crucially, Prof. Bui's mentorship and profound expertise in AI surrogates will provide the exact guidance I need to successfully integrate deterministic AC safety filters into learning-based policies, ultimately realizing my vision of a physically resilient smart grid.

\vspace{1ex}
\noindent \textbf{5. Career Objectives}

Following the completion of my Ph.D, I intend to deepen my expertise as an R\&D scientist within the U.S. energy sector before ultimately returning to Vietnam to contribute to the technical modernization of our national grid infrastructure. By repatriating advanced learning-based control methodologies, my long-term vision is to help transition our power architectures toward autonomous, resilient networks capable of stabilizing the rapid integration of clean energy sources. Ultimately, the Ph.D. program at the University of Michigan-Dearborn provides the exact interdisciplinary ecosystem and mentorship necessary to launch this career trajectory.

\end{document}
